% !TeX root = ../paper.tex
% 個人情報は docs/personal_data.md、実験情報は docs 内の資料に基づいて記入する。
\begingroup
\newcount\SubmissionHour
\SubmissionHour=\time
\divide\SubmissionHour by 60
\newcount\SubmissionMinute
\SubmissionMinute=\time
\advance\SubmissionMinute by -\numexpr60*\SubmissionHour\relax

\begin{center}
  {\huge{理工学基礎実験レポート}}
\end{center}

\vskip3\baselineskip

\begin{table}[!h]
\renewcommand{\arraystretch}{1.8}
 \centering
 \large
  \begin{tabular}{|p{25mm}|p{122mm}|}
   \hline
   実験日 &　年　月　日（火・金）（午前・午後） \\
   \hline
   実験題目 &   \\
   \hline
  \end{tabular}

\vskip1.5\baselineskip

  \begin{tabular}{|p{22.5mm}|p{38mm}|p{13mm}|p{13mm}|p{18mm}|p{28mm}|}
   \hline
   学科 &　& クラス &　& 学籍番号 &　\\
   \hline
   報告者氏名 & \multicolumn{5}{l|}{　} \\
   \hline
  \end{tabular}

\vskip1.5\baselineskip

  \begin{tabular}{|p{22.5mm}|p{60.5mm}|p{60.5mm}|}
   \hline
   　&  &　\\
   \cline{2-3}
   共同実験者 &　&　\\
   \cline{2-3}
   　&  &　\\
   \hline
  \end{tabular}

\vskip1.5\baselineskip

  \begin{tabular}{|p{39mm}|p{108mm}|}
   \hline
   レポート提出日 & \number\year 年\number\month 月\number\day 日\number\SubmissionHour 時\number\SubmissionMinute 分 \\
   \hline
   改訂レポート提出日 &　年　月　日　時　分 \\
   \hline
  \end{tabular}

\vskip1.5\baselineskip

\begin{flushleft}
  \begin{tabular}{|p{13.5mm}|p{60.5mm}|}
   \hline
   室温 &  \\
   \hline
   湿度 &  \\
   \hline
   気圧 &  \\
   \hline
  \end{tabular}
\end{flushleft}

\renewcommand{\arraystretch}{1.0}
\end{table}

\thispagestyle{empty}
\clearpage
\addtocounter{page}{-1}
\endgroup
